\documentclass[11pt,a4paper]{article}
\usepackage[T1]{fontenc}
\usepackage[utf8]{inputenc}
\usepackage{lmodern,microtype,amsmath,amssymb}
\usepackage[margin=21mm]{geometry}
\usepackage{booktabs,longtable,tabularx,array,enumitem,needspace,pdflscape}
\usepackage{xcolor,fancyhdr}
\usepackage[colorlinks=true,linkcolor=blue!45!black,urlcolor=blue!45!black,citecolor=blue!45!black]{hyperref}
\hypersetup{pdftitle={Experimental pp Rivet analyses in HerwigPolarizedPheno},pdfsubject={Catalogue of RHIC measurements and implemented analysis coverage}}
\setlength{\parindent}{0pt}
\setlength{\parskip}{5pt}
\setlength{\emergencystretch}{2em}
\setlength{\headheight}{14pt}
\setlist[description]{font=\normalfont\bfseries,leftmargin=0pt,labelindent=0pt,itemsep=3pt,parsep=0pt}
\pagestyle{fancy}
\fancyhf{}
\fancyhead[L]{\small Experimental pp Rivet analyses}
\fancyhead[R]{\small HerwigPolarizedPheno}
\fancyfoot[C]{\thepage}
\newcommand{\GeV}{\mathrm{GeV}}
\newcommand{\id}[1]{{\small\texttt{\detokenize{#1}}}}
\newcommand{\summaryname}[1]{\texttt{\detokenize{#1}}}
\newcommand{\sourcefile}[1]{\href{https://github.com/apapaefs/HerwigPolarizedPheno/blob/78d1be3962c3e058409634af28dc9cd52dc9b77d/analyses/rivet/pp/#1.cc}{\texttt{\detokenize{#1.cc}}}}
\newcommand{\repolink}[2]{\href{https://github.com/apapaefs/HerwigPolarizedPheno/blob/78d1be3962c3e058409634af28dc9cd52dc9b77d/#1}{#2}}
\begin{document}

\begin{landscape}
\thispagestyle{empty}
\begingroup
\begin{center}
{\Large\bfseries Experimental pp Rivet analyses: summary\par}
\vspace{4pt}
{\normalsize RHIC measurements implemented in HerwigPolarizedPheno}
\end{center}
\vspace{6pt}
\fontsize{10.0}{12.4}\selectfont
\setlength{\tabcolsep}{4pt}
\renewcommand{\arraystretch}{1.7}
\begin{tabularx}{\linewidth}{@{}l>{\raggedright\arraybackslash}Xlc>{\raggedright\arraybackslash}p{29mm}l@{}}
\toprule
\textbf{Rivet name} & \textbf{Experimental observable} & \textbf{arXiv paper} &
\shortstack{\textbf{Energy}\\$\sqrt{s}$ [GeV]} & \textbf{Beam type} & \textbf{Analysis source file}\\
\midrule
\summaryname{STAR_2019_I1708793} & $W^\pm$: $A_L(\eta_e)$, $A_{LL}(|\eta_e|)$; integrated $Z/\gamma^*$ $A_L$ & \href{https://arxiv.org/abs/1812.04817}{1812.04817} & 510 & Longitudinally polarized $pp$ & \sourcefile{STAR_2019_I1708793}\\
\summaryname{STAR_2021_I1850855} & Inclusive-jet $A_{LL}(p_T)$; same/opposite-sign dijet $A_{LL}(M_{jj})$ & \href{https://arxiv.org/abs/2103.05571}{2103.05571} & 200 & Longitudinally polarized $pp$ & \sourcefile{STAR_2021_I1850855}\\
\summaryname{STAR_2022_I1949588} & Inclusive-jet $A_{LL}(p_T)$; four dijet-topology $A_{LL}(M_{jj})$ measurements & \href{https://arxiv.org/abs/2110.11020}{2110.11020} & 510 & Longitudinally polarized $pp$ & \sourcefile{STAR_2022_I1949588}\\
\summaryname{PHENIX_2023_I2033856} & Inclusive/isolated prompt-photon invariant cross sections; isolated-photon $A_{LL}(p_T)$ & \href{https://arxiv.org/abs/2202.08158}{2202.08158} & 510 & Longitudinally polarized $pp$ & \sourcefile{PHENIX_2023_I2033856}\\
\bottomrule
\end{tabularx}
\vspace{10pt}
\normalsize
The collision energy is the proton--proton centre-of-mass energy. Equal-energy
beams have $E_p=100\,\GeV$ at $\sqrt{s}=200\,\GeV$ and
$E_p=255\,\GeV$ at $\sqrt{s}=510\,\GeV$.
The implemented predictions use separate $++,+-,-+,--$ physical-helicity samples.

\textbf{Coverage.} Four implementations correspond to four experimental publications:
three STAR analyses and one PHENIX analysis. The directory also contains three
internal Monte Carlo studies, which do not reproduce experimental data and are
excluded from this experimental table.

\textbf{Interpretation.} Every experimental \texttt{.info} file declares
\texttt{UNVALIDATED}. STAR weak bosons use a simplified dressed-lepton truth
selection; STAR jets use a generator-dependent hard-shower-parton proxy.
The PHENIX analysis is diagnostic because its current generated process omits
a fragmentation-photon event class. Details and reference-data coverage follow.

\textbf{Source snapshot.}
\href{https://github.com/apapaefs/HerwigPolarizedPheno/tree/78d1be3962c3e058409634af28dc9cd52dc9b77d/analyses/rivet/pp}{\texttt{analyses/rivet/pp/}},
commit \href{https://github.com/apapaefs/HerwigPolarizedPheno/commit/78d1be3962c3e058409634af28dc9cd52dc9b77d}{\texttt{78d1be3962c3}};
all 27 directory files checked against GitHub. Catalogue date: 30 September 2026.
\endgroup
\end{landscape}
\clearpage

\begin{center}
{\LARGE\bfseries Experimental proton--proton analyses\par}
{\large Rivet implementation catalogue\par}
30 September 2026
\end{center}

\section{Scope and interpretation}
This catalogue includes every experimental-measurement-linked Rivet analysis
implemented by a \texttt{.cc} file in
\href{https://github.com/apapaefs/HerwigPolarizedPheno/tree/78d1be3962c3e058409634af28dc9cd52dc9b77d/analyses/rivet/pp}{\texttt{analyses/rivet/pp/}}.
There are \textbf{four implementations: three STAR and one PHENIX}, each linked
to a distinct publication. All measurements concern longitudinally polarized
proton--proton collisions at RHIC. No experimental LHC pp implementation occurs
in this directory. The source snapshot is commit
\href{https://github.com/apapaefs/HerwigPolarizedPheno/commit/78d1be3962c3e058409634af28dc9cd52dc9b77d}{\texttt{78d1be3962c3}}.

The entries describe the implemented object definitions, selections, reference
inventory and coverage limits. Every experimental \texttt{.info} file declares
\texttt{Status: UNVALIDATED}. A verified reference transcription, a meaningful
generator-level comparison and a reproduction of detector response, triggers,
background subtraction and experimental corrections are separate requirements.
The present implementations do not establish the last of these.

\subsection*{Helicity observables and normalization}
Let $\sigma_{h_1h_2}(b)$ be the cross-section-normalized yield in bin $b$ from
independent, fully longitudinally polarized proton samples with
$h_1,h_2=\pm1$. The card labels \texttt{PP}, \texttt{PM}, \texttt{MP} and
\texttt{MM} denote $++,+-,-+,--$. They label helicity, whereas the
\texttt{p+} beam identifier denotes proton charge. Define
\begin{align}
 \sigma_{UU} &= \tfrac14(\sigma_{++}+\sigma_{+-}+\sigma_{-+}+\sigma_{--}),\\
 \Delta\sigma_{LL} &= \tfrac14(\sigma_{++}-\sigma_{+-}-\sigma_{-+}+\sigma_{--}),
 & A_{LL}&=\frac{\Delta\sigma_{LL}}{\sigma_{UU}},\\
 A_L^{(1)} &= \frac{\sigma_{++}+\sigma_{+-}-\sigma_{-+}-\sigma_{--}}
 {\sigma_{++}+\sigma_{+-}+\sigma_{-+}+\sigma_{--}},\\
 A_L^{(2)}&=\frac{\sigma_{++}-\sigma_{+-}+\sigma_{-+}-\sigma_{--}}
 {\sigma_{++}+\sigma_{+-}+\sigma_{-+}+\sigma_{--}}.
\end{align}
The workflow forms ratios after cross-section normalization of each physical
sample; raw histogram sums are not directly interchangeable across samples.
STAR weak-boson observables additionally use beam exchange and lepton
pseudorapidity reflection. Experimental polarization corrections belong to
the reference measurement; the ideal helicity samples do not introduce an
extra experimental beam-polarization factor into these ratios.

The current pp hard-process cards are \textbf{LO plus parton shower}.
NLO PDF inputs do not promote the hard calculation to NLO. The code fills
helicity yields; the campaign postprocessor constructs the published
asymmetries. Energies and momenta below are in GeV and angles are in radians.
\clearpage

\section{Complete inventory and shared jet definition}
\small
\renewcommand{\arraystretch}{1.2}
\begin{tabularx}{\textwidth}{@{}l>{\raggedright\arraybackslash}Xr@{}}
\toprule
Rivet identifier & Stored experimental reference coverage & Ref.\\
\midrule
\id{STAR_2019_I1708793} & Five series, 19 points: $W^\pm$ $A_L$, $W^\pm$ $A_{LL}$ and $Z/\gamma^*$ $A_L$ & \cite{r1708793}\\
\id{STAR_2021_I1850855} & Five series, 47 points; 36 in the primary covariance, with an 11-point alternative inclusive projection & \cite{r1850855}\\
\id{STAR_2022_I1949588} & Five series, 63 points; 59 in the primary covariance, excluding four low-$p_T$ inclusive bins & \cite{r1949588}\\
\id{PHENIX_2023_I2033856} & Three series, 43 points: two 18-point cross sections and seven isolated-photon $A_{LL}$ points & \cite{r2033856}\\
\bottomrule
\end{tabularx}
\normalsize
These are stored reference-series counts, not counts of statistically
independent measurements. In particular, the combined STAR 200 GeV inclusive
projection overlaps its central and forward projections.

\subsection*{STAR hard-shower-parton and stable-particle modes}
Both jet analyses use
\repolink{analyses/rivet/pp/STARPolarizedJets.hh}{\texttt{STARPolarizedJets.hh}}.
The default \texttt{LEVEL=HARDPARTON} follows terminal parton descendants of
outgoing signal-process-vertex lines and rejects beam-remnant descendants.
An unresolved hard-process provenance vetoes the event; a dedicated
\texttt{ProvenanceStatus} histogram records this diagnostic. The helper
clusters these partons using anti-$k_T$ with four-vector recombination.
The optional \texttt{LEVEL=HADRON} clusters stable visible final-state
particles and is a separate modelling diagnostic.

The default object is a \textbf{Herwig hard-shower-parton proxy}. Its outgoing-line
seeds retain final-state radiation but omit initial-state-radiation emissions.
STAR's corrected parton-jet target includes initial- and final-state radiation
and uses detector embedding, jet matching and trigger/reconstruction
corrections tied to the experiment's generator model. Applying STAR detector
bin intervals directly to Herwig parton jets, then plotting at corrected
representative parton coordinates, does not reproduce that correction chain.

The nominal jet cards use MPI off and a 4 GeV generator cut. The registered
3/4/5 GeV stability scan addresses the generator cut; it does not validate the
experimental object definition. Hadron/MPI-on and shower-spin-off families
remain diagnostics. The
\repolink{docs/experimental-analysis-compatibility-audit.md}{experimental-analysis compatibility audit}
records these distinctions and the weak-boson/photon coverage limits.

\subsection*{Internal studies excluded from the experimental catalogue}
\id{MC_POLDIJETS}, \id{MC_POLJETSHAPES} and
\id{MC_POLJETSHAPES_HARD} are internal shower-spin studies, without
corresponding experimental data series. Literature references in their metadata
do not make them experimental implementations. Their files remain part of the
27-file source-inventory check.
\clearpage

\section{STAR weak-boson analysis}
\subsection{\texorpdfstring{\id{STAR_2019_I1708793}}{STAR 2019 I1708793}}
\begin{description}
\item[Paper and source.]
STAR's longitudinal single- and double-spin asymmetries in $W$ and
$Z/\gamma^*$ production~\cite{r1708793};
\sourcefile{STAR_2019_I1708793}.
\item[Beams and energy.]
Longitudinally polarized $pp$, $\sqrt{s}=510\,\GeV$;
two $255\,\GeV$ proton beams, with separate four-helicity samples for the
$W^+$, $W^-$ and $Z/\gamma^*$ hard channels.
\item[Implemented objects and selections.]
Electrons/positrons are dressed with photons within $\Delta R=0.1$, using
Rivet's \texttt{NODECAY} lepton and photon origin flags. The $W$ selection
requires exactly one selected dressed lepton with
$25<p_T^e<50\,\GeV$ and $|\eta_e|<1.5$; its charge selects $W^+$ or $W^-$.
These are \emph{transverse-momentum} cuts in the \texttt{.cc} source, although
the \texttt{.info} description calls them transverse-energy cuts.
There is no explicit boson-ancestry requirement.

The $Z/\gamma^*$ selection counts opposite-charge dressed pairs with
$p_T^e>14\,\GeV$, $|\eta_e|<1.1$ and
$70<m_{ee}<110\,\GeV$. Each lepton must carry more than 88\% of the total
transverse energy in a $\Delta R<0.7$ cone, including its own energy and
excluding its dressing constituents and neutrinos from the added cone sum.
\item[Observables.]
The code fills six $0.5$-wide signed-$\eta_e$ bins per $W$ charge and one
integrated $Z/\gamma^*$ yield. Postprocessing forms
$A_L^W(\eta)=\tfrac12[A_L^{(1)}(\eta)+A_L^{(2)}(-\eta)]$ and folds the
four-helicity $W$ yields into three $|\eta_e|$ bins for $A_{LL}$.
In the symmetric integrated $Z/\gamma^*$ acceptance, the mean single-spin
asymmetry is $(\sigma_{++}-\sigma_{--})/\sum\sigma$.
\item[Reference data.]
The
\href{https://doi.org/10.17182/hepdata.105912.v1}{HEPData record 105912},
Tables 17 and 18, supplies six combined $W^+$ and six combined $W^-$
$A_L$ points. Publication Table I supplies three $A_{LL}$ points per charge;
the paper supplies one integrated $Z/\gamma^*$ $A_L$ point. The normalized
\repolink{data/phenomenology/STAR_2019_I1708793/reference.json}{reference snapshot}
therefore contains five series and 19 points. The compressed reference YODA
contains all five physics series, together with statistical-only companion
objects. Its assigned contiguous theory-$\eta$ edges are not literal STAR
detector-acceptance boundaries; the HEPData tables give effective point coordinates.
\item[Coverage and limitations.]
This is a simplified LO-plus-shower truth comparison. The $W$ selection lacks
the experiment's local/cone isolation, momentum balance, opposite-side energy,
trigger, tracking, purity and background treatment; its symmetric six-bin
acceptance also differs from STAR's barrel and one-sided endcap coverage.
The published combined covariance is unavailable. Plotted $W$ $A_L$ point
systematic errors omit correlated content, including a 3.3\% polarization
scale; the paper-only $A_{LL}$ series instead retain publication totals.
Diagonal residuals do not define a complete experimental goodness of fit.
The metadata status is \texttt{UNVALIDATED}.
\end{description}
\clearpage

\section{STAR jet and dijet analyses}
\subsection{\texorpdfstring{\id{STAR_2021_I1850855}}{STAR 2021 I1850855}}
\begin{description}
\item[Paper and source.]
STAR's inclusive-jet and dijet longitudinal double-spin asymmetries at
$\sqrt{s}=200\,\GeV$~\cite{r1850855};
\sourcefile{STAR_2021_I1850855}.
\item[Beams and objects.]
Longitudinally polarized $pp$, two $100\,\GeV$ beams.
Anti-$k_T$ jets with $R=0.6$, default \texttt{LEVEL=HARDPARTON}; the shared
object and its limitations are described in Section 2.
\item[Inclusive selection.]
At most the two highest-$p_T$ jets passing $p_T\geq6\,\GeV$ and
$|\eta|<1$ contribute. Central jets have $|\eta|<0.5$ and forward jets
have $0.5\leq|\eta|<1$; an additional combined $|\eta|<1$ projection
is retained. Each projection has 11 source $p_T$ bins spanning 6--37.3 GeV.
\item[Dijet selection.]
The two highest-$p_T$ jets in the clustered list must satisfy
$p_{T,1}>8\,\GeV$, $p_{T,2}>6\,\GeV$,
$|\eta_1|,|\eta_2|<0.8$ and $\Delta\phi>2\pi/3$.
The source separates same-sign ($\eta_1\eta_2>0$) and opposite-sign
($\eta_1\eta_2<0$) pairs; a zero product enters neither category.
Each dijet invariant-mass histogram has seven bins spanning 17--82 GeV.
\item[Observables and reference data.]
The five yield histograms produce $A_{LL}(p_T)$ for three inclusive
projections and $A_{LL}(M_{jj})$ for two dijet categories. The
\href{https://doi.org/10.17182/hepdata.104836.v1}{HEPData record 104836}
contains 21 tables: supporting yield/bin information in Tables 1--3,
asymmetry data in Tables 4--8 ($3\times11+2\times7=47$ points), and
correlations in Tables 9--21. The
\repolink{data/phenomenology/STAR_2021_I1850855/reference.json}{normalized snapshot}
stores all five series. The compressed reference YODA also contains all five
series, including the alternative combined inclusive projection.
\item[Covariance and coverage.]
The primary set uses the central and forward inclusive projections and both
dijet categories, giving a 36-point covariance. The combined 11-point
inclusive projection is an alternative display of overlapping data and is
excluded from this simultaneous fit. Published point-to-point covariance
blocks are combined with separately profiled global uncertainties: an
absolute $7.0\times10^{-4}$ relative-luminosity shift and a 6.1\%
polarization scale.

The source-level cuts are implemented on Herwig shower-parton jets, using
detector-bin intervals with corrected parton display coordinates. This remains
a generator-dependent LO-plus-shower modelling comparison, rather than a
reproduction of STAR's detector selection and correction chain. The optional
stable-particle mode does not remove that distinction. Metadata:
\texttt{UNVALIDATED}.
\end{description}
\clearpage

\subsection{\texorpdfstring{\id{STAR_2022_I1949588}}{STAR 2022 I1949588}}
\begin{description}
\item[Paper and source.]
STAR's inclusive-jet and dijet longitudinal double-spin asymmetries at
$\sqrt{s}=510\,\GeV$~\cite{r1949588};
\sourcefile{STAR_2022_I1949588}.
\item[Beams and objects.]
Longitudinally polarized $pp$, two $255\,\GeV$ beams;
anti-$k_T$ with $R=0.5$, default \texttt{LEVEL=HARDPARTON}.
\item[Inclusive selection.]
At most the two highest-$p_T$ accepted jets contribute, with
$p_T\geq7\,\GeV$ and $|\eta|<0.9$. Fourteen source bins span 7--62.8 GeV.
The registry treats the first four inclusive bins as diagnostics; the primary
comparison begins at the source lower edge $p_T=13.1\,\GeV$. This restriction
does not change the \texttt{.cc} histogram definition.
\item[Dijet selection and topology.]
The two highest-$p_T$ jets must have $p_{T,1}>7\,\GeV$,
$p_{T,2}>5\,\GeV$, $\Delta\phi>2\pi/3$ and
$|\eta_1-\eta_2|<1.6$. The four topology categories are:
\end{description}
\small
\renewcommand{\arraystretch}{1.15}
\begin{tabularx}{\textwidth}{@{}c>{\raggedright\arraybackslash}Xrc@{}}
\toprule
Topology & Source pseudorapidity selection & Points & Source $M_{jj}$ range [GeV]\\
\midrule
A & Both $0.3<|\eta|<0.9$, same-sign $\eta_1\eta_2>0$ & 12 & 12--101\\
B & One $|\eta|<0.3$, one $0.3<|\eta|<0.9$ & 13 & 12--121\\
C & Both $|\eta|<0.3$ & 12 & 12--101\\
D & Both $0.3<|\eta|<0.9$, opposite-sign & 12 & 14--121\\
\bottomrule
\end{tabularx}
\normalsize
\begin{description}
\item[Observables and reference data.]
Inclusive $A_{LL}(p_T)$ and four $A_{LL}(M_{jj})$ series.
\href{https://doi.org/10.17182/hepdata.114778.v1}{HEPData record 114778}
provides 14 inclusive and $12+13+12+12=49$ dijet points: 63 in total.
Tables 6--20 supply covariance information; the
\repolink{data/phenomenology/STAR_2022_I1949588/reference.json}{normalized snapshot}
and compressed reference YODA contain all five series. The primary covariance is the 59-point submatrix
remaining after the four low-$p_T$ inclusive bins are excluded.
\item[Covariance and coverage.]
The reported systematic already includes the absolute $4.7\times10^{-4}$
relative-luminosity term, whereas the published correlation tables exclude
it. The current reference separates the point-to-point component as
$\sqrt{\max(\delta_{\mathrm{sys,tot}}^2-(4.7\times10^{-4})^2,0)}$
and profiles the luminosity shift once. The polarization scale is 6.4\%.

As for the 200 GeV implementation, the hard-shower-parton object and direct
application of detector-bin intervals remain generator-dependent proxies.
Generator-cut stability and corrected covariance require separately verified
campaign/postprocessing provenance; these source definitions do not certify
old production results. Metadata: \texttt{UNVALIDATED}.
\end{description}
\clearpage

\section{PHENIX prompt-photon analysis}
\subsection{\texorpdfstring{\id{PHENIX_2023_I2033856}}{PHENIX 2023 I2033856}}
\begin{description}
\item[Paper and source.]
PHENIX's midrapidity direct-photon cross sections and longitudinal double-spin
asymmetry at $\sqrt{s}=510\,\GeV$~\cite{r2033856};
\sourcefile{PHENIX_2023_I2033856}.
\item[Beams and photon selection.]
Longitudinally polarized $pp$, two $255\,\GeV$ beams. The analysis counts
every Rivet prompt photon satisfying $|\eta_\gamma|<0.25$ and
$6<p_T^\gamma<30\,\GeV$, including non-decay photons produced by the shower.
The inclusive and isolated cross-section histograms each have 18 bins over
6--30 GeV. The isolated-photon helicity-yield histogram has seven bins over
6--20 GeV.
\item[Isolation.]
Within $\Delta R<0.5$, the code sums the \emph{energy} of other final-state
particles, excluding the candidate photon and neutrinos. Charged particles
enter for $p_T\geq0.2\,\GeV$ and neutral particles for $E\geq0.3\,\GeV$.
Isolation requires $E_{\mathrm{cone}}<0.1E_\gamma$.
The charged threshold is on transverse momentum in this implementation;
the publication describes a total-momentum threshold of 200 MeV/$c$.
The code applies full-azimuth truth acceptance, with no detector response.
\item[Observables and normalization.]
The two cross sections are the inclusive and isolated invariant distributions
$E\,d^3\sigma/d^3p$, in $\mathrm{pb}\,\GeV^{-2}$.
Each photon contributes $1/p_T$ to the corresponding histogram and the
cross-section normalization includes $1/(2\pi\Delta\eta)$ with
$\Delta\eta=0.5$; bin-width normalization gives
\begin{equation}
 E\frac{d^3\sigma}{d^3p}
 =\frac{1}{2\pi p_T\Delta\eta}\frac{d\sigma}{dp_T}.
\end{equation}
The isolated-photon $A_{LL}$ is constructed from ordinary isolated-photon
helicity yields, without the $1/p_T$ cross-section weighting.
\item[Reference data.]
\href{https://doi.org/10.17182/hepdata.129088.v1}{HEPData record 129088},
Tables 1 and 2, supplies two 18-point cross-section series and one seven-point
$A_{LL}$ series: three series and 43 points in the
\repolink{data/phenomenology/PHENIX_2023_I2033856/reference.json}{normalized snapshot}.
The compressed reference YODA contains all three series. Global uncertainties
in the stored HEPData-based reference are 10\% cross-section luminosity,
an absolute $3.9\times10^{-4}$ asymmetry shift and a 6.6\% polarization scale.
The arXiv v2 body quotes $3.8\times10^{-4}$ for the shift; the snapshot follows
the $3.9\times10^{-4}$ value supplied by HEPData.
\item[Coverage and limitations.]
The current cards use a direct $\gamma+$jet hard matrix element with a QED
shower. They omit the separately normalized QCD dijet event class that can
produce a fragmentation/shower photon. Isolation reduces that contribution
but does not eliminate it. The inclusive prediction is consequently
incomplete; the isolated cross section and $A_{LL}$ also remain conditional.
A process-complete prediction needs an additional, separately normalized
fragmentation-photon sample with direct/fragmentation overlap removal, or a
validated NLO prompt-photon implementation. The compatibility
audit therefore retains this entry as \textbf{diagnostic only}, excluded from
publication-grade conclusions. Metadata: \texttt{UNVALIDATED}.
\end{description}
\clearpage

\section*{References and reproducibility}
Each entry links the implementation at the frozen repository commit; its
matching \texttt{.info} and \texttt{.plot} files are in the same directory.
Normalized reference snapshots are stored in
\texttt{data/phenomenology/\{analysis\}/reference.json}; source manifests and
raw HEPData tables record the transcription provenance. The matching registries
in \texttt{config/phenomenology/} specify cards, sample families and, where
available, covariance and global uncertainties. The compressed reference YODA
contains the physics series described here; statistical-only companions are
not additional measurements. Complete covariance, nuisance treatment and
provenance still require the normalized snapshots and registries.

The catalogue inventories implemented analysis definitions and experimental
references, rather than assessing a particular generated campaign. No
generator accuracy, experimental response equivalence or complete goodness
of fit follows from the presence of an analysis file or a reference overlay.

\begin{thebibliography}{4}
\bibitem{r1708793} STAR Collaboration,
\emph{Measurement of the longitudinal spin asymmetries for weak boson production in proton-proton collisions at $\sqrt{s}=510$ GeV},
Phys. Rev. D 99 (2019) 051102.
\href{https://doi.org/10.1103/PhysRevD.99.051102}{doi:10.1103/PhysRevD.99.051102};
\href{https://arxiv.org/abs/1812.04817}{arXiv:1812.04817}.

\bibitem{r1850855} STAR Collaboration,
\emph{Longitudinal double-spin asymmetry for inclusive jet and dijet production in polarized proton collisions at $\sqrt{s}=200$ GeV},
Phys. Rev. D 103 (2021) L091103.
\href{https://doi.org/10.1103/PhysRevD.103.L091103}{doi:10.1103/PhysRevD.103.L091103};
\href{https://arxiv.org/abs/2103.05571}{arXiv:2103.05571}.

\bibitem{r1949588} STAR Collaboration,
\emph{Longitudinal double-spin asymmetry for inclusive jet and dijet production in polarized proton collisions at $\sqrt{s}=510$ GeV},
Phys. Rev. D 105 (2022) 092011.
\href{https://doi.org/10.1103/PhysRevD.105.092011}{doi:10.1103/PhysRevD.105.092011};
\href{https://arxiv.org/abs/2110.11020}{arXiv:2110.11020}.

\bibitem{r2033856} PHENIX Collaboration,
\emph{Measurement of Direct-Photon Cross Section and Double-Helicity Asymmetry at $\sqrt{s}=510$ GeV in $\vec p+\vec p$ Collisions},
Phys. Rev. Lett. 130 (2023) 251901.
\href{https://doi.org/10.1103/PhysRevLett.130.251901}{doi:10.1103/PhysRevLett.130.251901};
\href{https://arxiv.org/abs/2202.08158}{arXiv:2202.08158}.
\end{thebibliography}
\end{document}
